\subsection{Résultats, Analyse et Discussion}

\subsubsection{Determination de la celerite}

Nous mesurons pour des phases a 0 $\pm 5 ^\circ$ sur en Annexe (voir Tab. \ref{tab:lambda}) :

Dont nous pouvons extraire une moyenne de 0.83 $cm$ quand on considere la longueur d'onde de 1.8 $cm$ comme la somme de deux intervalles de phase nulle de nos signaux. En effet, si on le divise par deux, nous obtenons 0.9 $cm$ qui est relativement proche de la moyenne.  \\\\
Grace a la relation $\lambda$ = c $\cdot$ $T$ avec T = $\frac{1}{f}$ (voir Eq. \ref{eq:lambda}), nous obtenons : 
\[
	\mathrm{c} = \lambda \cdot f = 0.0083 \cdot 40'000 = 3 32 \pm 4 m/s
\]


En supposant que la fréquence $f$ est connue avec une précision négligeable.
L'incertitude relative de $c$ est donc donnée par :
\[
\frac{\Delta c}{c} = \frac{\Delta \lambda}{\lambda}.
\]

L'incertitude sur la longueur d'onde est celle de la règle utilisée soit
\[
\Delta \lambda = \pm 0{,}1~\mathrm{cm} = 1 \cdot 10^{-4}~\mathrm{m},
\]
alors l'incertitude absolue sur $c$ est
\[
\Delta c = c \cdot \frac{\Delta \lambda}{\lambda}
= 332 \cdot \frac{1 \cdot 10^{-4}}{8,3 \cdot 10^{-3}}
= 4 ~\mathrm{m/s}.
\]

\[
	c = (3,3 \pm 0,4) \cdot 10^{2}~\mathrm{m/s}\qquad pour~ c_t_h \approx 343 \mathrm{m/s}.
\]
\subsubsection{Nature spherique des ondes accoustiques}

\begin{figure}[h!]
    \centering
    \includegraphics[width=0.5\textwidth]{accoustique/angle.png}
    \caption{Graphique polaire $U(\theta)$}
    \label{fig:angle}
\end{figure}

\noindent
Le graphique (Fig. \ref{fig:angle}) montre que la tension mesurée (Voir Tab. \ref{tab:angle}) est maximale pour $\theta = 0^\circ$ et décroît lorsque l’angle augmente. EXPLIQUER LIEN AVEC ONDES SPHERIQUES



\subsubsection{Vérification de la variation de la pression}
La Fig.~\ref{fig:U(1/R)} présente l’évolution de la tension mesurée en fonction de l’inverse de la distance entre l’émetteur et le récepteur.
La tension délivrée par le capteur piézoélectrique étant proportionnelle à la pression acoustique incidente, cette représentation permet de visualiser directement la variation de la pression avec la distance de propagation de l’onde.

\begin{center}
	\includegraphics[scale=0.4]{accoustique/acc1.png}
\captionof{figure}{Variation de la tension en fonction de l’inverse de la distance}
\label{fig:U(1/R)}
\end{center}

Afin de vérifier le modèle théorique de décroissance de la pression acoustique, nous avons ensuite tracé le carré de la tension en fonction de l’inverse du carré de la distance (Fig.~\ref{fig:U2(1/R2)}). Nous verifions que la linearité reste inchangée bien que l'offset bouge.
Nous constatons egalement sur Fig. \ref{fig:U(1/R)inv}, que si nous inversons le branchement E-R, la pente devient beaucoup certe plus petite, mais que la linéarité est conservée.

\begin{figure}[h!]
    \centering
    % Première figure
    \begin{minipage}[b]{0.48\linewidth}
        \centering
        \includegraphics[width=\linewidth]{accoustique/acc2.png}
	    \caption{U$^2$(1/D$^2$)}
        \label{fig:U2(1/R2)}
    \end{minipage}
    \hfill
    % Deuxième figure
    \begin{minipage}[b]{0.48\linewidth}
        \centering
        \includegraphics[width=\linewidth]{accoustique/acc3.png}
	    \caption{U(1/D) avec E-R inversés}
        \label{fig:U(1/R)inv}
    \end{minipage}
\end{figure}

En somme, nous observons dans les 3 cas une relation lineaire avec les pente positives avec similaires incertitudes telles que : 

\[
	$\alpha_U = 155 \pm 8$ V/m ,\qquad
	$\alpha_U_^2$ $= 30067 \pm 1771$ $V^2$/$m^2$,
	\qquad $\alpha_U_i_n_v = 0.21 \pm 0.02$ $V/m$.
\]


